Astronomy has benefited enormously from the constant improvements in computing power and the advances in fields like data mining and artificial intelligence. Specifically, modern hydrodynamic simulations have provided massive data sets of galaxies with unprecedented accuracy.

Leveraging these advancements, the work presented here presents the results of applying clustering techniques in the dynamic decomposition of galaxies, a field that seeks an understanding of how different parts of a galaxy contribute to its movements and general evolution. We compare the results obtained with the ones generated by already established methods in the area. Specifically, we will evaluate Hierarchical Clustering, Fuzzy C-Means, and Ensemble Agglomerative Clustering, together with two outliers removal techniques: one from coming from astronomy (Radial distance cut) and another one coming from data mining (Isolation Forest).

The results show that both Hierarchical Clustering and Fuzzy C-Means can get similar results to the already established methods used in the field, although with significant limitations, especially in identifying which component corresponds to which group. On the other hand, we found that internal validation metrics used on conventional clustering aren't useful for analyzing the performance of these techniques due to the nature of astronomy data. As to Evidence Accumulation Clustering, it displayed results that are far from the expected ones and at a very high computational cost, which discourages its usage in the way we explored it. An important discovery was the inability to determine if removing outliers benefits the clustering process in this specific context.

This thesis opens a path for future research in the field, including the extension of the dataset to include spiral galaxies, the exploration of normalization techniques and additional preprocess work, and the evaluation of other outliers removal methods and features of the data set that can be relevant for the dynamical decomposition of galaxies.

\vfill

\textbf{Keywords:} Data Minning - Unsupervised learning - Clustering - Outlier removal - Numerical simulations - Dynamic decomposition - Circularity parameter